\subsection{导子的一般定义}

考虑一个交换环 K，六个类型为 \(\Delta\) 的分次 K-模：A, A', A'', B, B', B''，以及三个 K-线性映射

\[
\mu \colon \mathrm{A} \times \mathrm{A} ^ {\prime} \rightarrow \mathrm{A} ^ {\prime \prime}, \quad \lambda_ {1} \colon \mathrm{B} \times \mathrm{A} ^ {\prime} \rightarrow \mathrm{B} ^ {\prime \prime}, \quad \lambda_ {2} \colon \mathrm{A} \times \mathrm{B} ^ {\prime} \rightarrow \mathrm{B} ^ {\prime \prime}
\]

使得相应的 K-线性映射

\[
\mathrm{A} \otimes_ {\mathrm{K}} \mathrm{A} ^ {\prime} \rightarrow \mathrm{A} ^ {\prime \prime}, \quad \mathrm{B} \otimes_ {\mathrm{K}} \mathrm{A} ^ {\prime} \rightarrow \mathrm{B} ^ {\prime \prime}, \quad \mathrm{A} \otimes_ {\mathrm{K}} \mathrm{B} ^ {\prime} \rightarrow \mathrm{B} ^ {\prime \prime}
\]

是次数为 0 的分次映射。像 \(\mu(a, a')\) （对于 \(a \in A\) ， \(a' \in A'\) ）简记为 \(a, a'\) 甚至 \(aa'\) ，另外两个双线性映射同理。因此 \(a, a'\) 的次数是 \(a\) 与 \(a'\) 的次数之和。

定义 1. 在上述设定和一个交换因子 \(\varepsilon\) （ \(\Delta \times \Delta\) 上的）之下，一个 \(\varepsilon\) -导子（或 \((\mathbf{K}, \varepsilon)\) -导子）（次数为 \(\delta \in \Delta\) ，从 \((\mathbf{A}, \mathbf{A}', \mathbf{A}'')\) 到 \((\mathbf{B}, \mathbf{B}', \mathbf{B}'')\) ）指的是次数为 \(\delta\) 的分次 K-线性映射的一个三元组：

\[
d \colon \mathrm{A} \rightarrow \mathrm{B}, \quad d ^ {\prime} \colon \mathrm{A} ^ {\prime} \rightarrow \mathrm{B} ^ {\prime}, \quad d ^ {\prime \prime} \colon \mathrm{A} ^ {\prime \prime} \rightarrow \mathrm{B} ^ {\prime \prime}
\]

的分次K-线性映射三元组，使得对每个齐次元素 \(a \in A\) 以及每个元素 \(a' \in A'\)

\[
d ^ {\prime \prime} (a. a ^ {\prime}) = (d a). a ^ {\prime} + \varepsilon_ {\delta , \deg (a)} a. (d ^ {\prime} a ^ {\prime}).\tag{4}
\]

显然，由线性性，只需在 \(a\) 和 \(a'\) 分别遍历 A 和 A' 的生成系时验证关系 (4) 即可。

把这三个映射 \(d, d', d''\) 用同一个字母 \(d\) 来表示往往是方便的（其合理性在于同样用 \(d\) 表示次数为 \(\delta\) 的分次 K-线性映射

\[
(a, a ^ {\prime}, a ^ {\prime \prime}) \mapsto (d a, d ^ {\prime} a ^ {\prime}, d ^ {\prime \prime} a ^ {\prime \prime})
\]

它把 \(\mathbf{A} \oplus \mathbf{A}' \oplus \mathbf{A}''\) 映到 \(\mathbf{B} \oplus \mathbf{B}' \oplus \mathbf{B}''\) 中。于是关系 (4) 可以更简单地写成

\[
d (a. a ^ {\prime}) = (d a). a ^ {\prime} + \varepsilon_ {\delta . \deg (a)} a. (d a ^ {\prime}).\tag{5}
\]

从 (A, A', A'') 到 (B, B', B'') 的给定次数的 \(\varepsilon\) -导子构成分次线性映射 K-模的一个子 K-模

\[
\operatorname{Homgr} _ {\mathbb {K}} (\mathrm{A} \oplus \mathrm{A} ^ {\prime} \oplus \mathrm{A} ^ {\prime \prime}, \mathrm{B} \oplus \mathrm{B} ^ {\prime} \oplus \mathrm{B} ^ {\prime \prime}).
\]

当 \(\varepsilon(\alpha, \beta) = 1\) 对所有 \(\alpha, \beta\) （ \(\Delta\) 中的）成立时，我们简称导子（或 K-导子）而不称 \(\varepsilon\) -导子。导子构成下列 K-模的一个子 K-模

\[
\operatorname{Hom} _ {\mathbb {K}} (\mathrm{A} \oplus \mathrm{A} ^ {\prime} \oplus \mathrm{A} ^ {\prime \prime}, \mathrm{B} \oplus \mathrm{B} ^ {\prime} \oplus \mathrm{B} ^ {\prime \prime}).
\]

当 \(\Delta = Z\) 且 \(\varepsilon(p, q) = (-1)^{pq}\) 时，每个偶次数的 \(\varepsilon\) -导子都是导子；每个奇次数的 \(\varepsilon\) -导子通常称为反导子（或 K-反导子）；于是反导子 d 满足

\[
d (a. a ^ {\prime}) = (d a). a ^ {\prime} + (- 1) ^ {\deg (a)} a. (d a ^ {\prime})\tag{6}
\]

对齐次元素 \(a \in A\) 成立。

注记。(1) 导子的概念可以推广到非分次模上，只需约定赋予这些模平凡分次结构即可。

\begin{enumerate}
\def\labelenumi{(\arabic{enumi})}
\setcounter{enumi}{1}
\tightlist
\item
  如果只考虑 \(\varepsilon\) -导子中次数为给定值 \(\delta\) 的那些，交换因子 \(\varepsilon\) 可以按如下方式消去：把双线性映射 \(\lambda_2: \mathbf{A} \times \mathbf{B}' \to \mathbf{B}''\) 替换为双线性映射 \(\lambda_2': \mathbf{A} \times \mathbf{B}' \to \mathbf{B}''\) ，使得对每个齐次元素 \(a\) （ \(\mathbf{A}\) 中的）以及所有 \(b' \in \mathbf{B}'\) ，
\end{enumerate}

\[
\lambda_ {2} ^ {\prime} (a, b ^ {\prime}) = \varepsilon_ {\delta , \deg (a)} \lambda_ {2} (a, b ^ {\prime}).
\]

那么 \(d\) 是关于双线性映射 \(\mu, \lambda_{1}, \lambda_{2}^{\prime}\) 的导子。

上述给出的 \(\varepsilon\) -导子的一般定义特别用于两种情况：

情形（I）：A = B， \(A' = B'\) , \(A'' = B''\) 且三个双线性映射 \(\mu\) , \(\lambda_{1}\) , \(\lambda_{2}\) 等于同一个映射。

情形（II）： \(\mathbf{A} = \mathbf{A}' = \mathbf{A}''\) ， \(\mathbf{B} = \mathbf{B}' = \mathbf{B}''\) ，因此（对于 \(\mu\) ） \(\mathbf{A}\) 是一个分次代数，且这两个 K-双线性映射

\[
\lambda_ {1} \colon \mathrm{B} \times \mathrm{A} \rightarrow \mathrm{B}, \quad \lambda_ {2} \colon \mathrm{A} \times \mathrm{B} \rightarrow \mathrm{B}\tag{7}
\]

使得相应的 K-线性映射 \(B \otimes_{K} A \to B\) ， \(A \otimes_{K} B \to B\) 是次数为 0 的分次映射。此时，一个 \(\varepsilon\) -导子（次数为 \(\delta\) ，从 A 到 B）是指一个分次 K-线性映射 \(d: A \to B\) （次数为 \(\delta\) ），使得对 A 中每个齐次元素 x 以及每个 \(y \in A\) ，我们有关系式

\[
d (x y) = (d x) y + \varepsilon_ {\delta , \deg (a)} x (d y).\tag{8}
\]

特别地，在情形 (II) 中考虑这样的情形：A 是带单位元的结合 K-代数，且 \(\lambda_{1}\) 和 \(\lambda_{2}\) 是一个 (A, A)-双模（§ 4，第 3 号，定义 3）的外部法则。特别当 A 和 B 是两个带单位元的结合 K-代数、给定分次 K-代数的一个带单位元的同态 \(\rho: A \to B\) 并在 B 上考虑由这两个外部法则定义的 (A, A)-双模结构时，上述情形成立

\[
\lambda_ {1} \colon (b, a) \mapsto b \rho (a), \quad \lambda_ {2} \colon (a, b) \mapsto \rho (a) b
\]

对于 \(a\in \mathbf{A},b\in \mathbf{B}.\)

情形（I）和（II）具有以下共同情形：考虑一个分次 K-代数 A，取 B = A，映射 (7) 均为 A 上的乘法。此时我们谈论分次代数 A 的 ε-导子（或 (K, ε)-导子）：它是 A 到自身的、次数为 δ 的分次 K-线性映射，对 A 中每个齐次元素 x 以及所有 \(y \in A\) 满足 (8)。特别地，若 A 是一个分次环（视为（结合的）Z-代数），我们就谈论环 A 的 ε-导子。

设A是一个含单位元的交换结合K-代数，B是一个A-模；当我们谈到从A到B的导子时，总是指B上由A-模结构导出的A-双模结构；此时公式

\[
d (x y) = x (d y) + y (d x) \quad \text { for } \quad x \in \mathrm{A}, y \in \mathrm{A}\tag{9}
\]

对这样的导子 \(d: A \rightarrow B\) 成立。

